\chapter*{Course Map}
\phantomsection%
\addcontentsline{toc}{chapter}{Course Map}
\markboth{Course Map}{}

Data science turns questions and data into evidence that others can inspect, reproduce, and challenge. This course emphasises analytical workflow, with R as a tool. Four stages build on Chapter~\ref{ch:toolkit}, which introduces R, RStudio, optional conda management, and the required Git/GitHub SSH workflow; the capstone integrates them. The chapters use varied datasets to build transferable skills rather than familiarity with a single dataset.

\begin{utpfigure}
\newcommand{\coursemapcard}[3]{%
  \parbox[c][1.52cm][t]{3.35cm}{\raggedright
    {\scriptsize\bfseries\color{UTPTeal}\strut #1}\par
    \vspace{0.02em}%
    {\bfseries\strut #2}\par
    \vspace{0.08em}%
    {\scriptsize\color{UTPMuted!80!UTPNavy}\strut #3\strut}%
  }%
}
\begin{tikzpicture}[
  node distance=0.58cm and 0.66cm,
  card/.style={draw=UTPBlue,fill=UTPPaleBlue,rounded corners=2mm,
    minimum width=3.99cm,minimum height=1.88cm,align=left,
    inner xsep=3.2mm,inner ysep=1.6mm,font=\small\color{UTPNavy}},
  finalcard/.style={card,draw=UTPTeal,fill=UTPPaleTeal},
  arr/.style={-{Stealth[length=2.25mm,width=1.55mm]},line width=0.9pt,
    draw=UTPTeal,shorten >=0.8pt,shorten <=0.8pt}
]
\node[card] (a) {\coursemapcard{STAGE 01}{Foundations}{Questions, data quality, and reproducibility}};
\node[card,right=of a] (b) {\coursemapcard{STAGE 02}{Working with data}{Wrangling, EDA, and visualization}};
\node[card,right=of b] (c) {\coursemapcard{STAGE 03}{Prediction}{Inference, regression, and classification}};
\node[card,below=of c] (d) {\coursemapcard{STAGE 04}{Structure}{Clustering and PCA}};
\node[finalcard,left=of d] (e) {\coursemapcard{STAGE 05}{Capstone}{An end-to-end analysis}};
\draw[arr] (a) -- (b);
\draw[arr] (b) -- (c);
\draw[arr] (c) -- (d);
\draw[arr] (d) -- (e);
\end{tikzpicture}
\end{utpfigure}

\begin{utptable}
\begin{tabularx}{\linewidth}{@{}>{\bfseries\leavevmode\color{UTPNavy}\centering\arraybackslash}m{0.08\linewidth}m{0.31\linewidth}X@{}}
\toprule
\rowcolor{UTPPaleBlue!92}
\utpth{Stage} & \utpth{Focus} & \utpth{Key question} \\
\midrule
1 & Workflow and data quality & Is the question answerable, and are the data fit for it? \\
2 & Wrangling, EDA, and visualization & What patterns do the data show, and how are they communicated? \\
3 & Inference, prediction, and validation & How certain is an estimate, and how well does a model generalise? \\
4 & Clustering and PCA & What structure exists when there is no label to predict? \\
5 & Capstone project & Can the whole workflow be defended end to end? \\
\bottomrule
\end{tabularx}
\end{utptable}

\section*{The workflow and the labs}
Every lab follows the same workflow, revisited at increasing depth, and works on a different dataset.

\begin{utpfig}[Course workflow]{The course revisits the same workflow at increasing depth. Each lab sits immediately after the chapter that supplies its concepts.}{fig:course-workflow}
\begin{tikzpicture}[x=\linewidth,y=1cm,
  stage/.style={draw=UTPBlue!70,fill=UTPPaleBlue,rounded corners=1.5mm,
    line width=0.7pt,text width=0.155\linewidth,minimum height=1.65cm,
    inner xsep=1.8mm,inner ysep=2mm,align=center,
    font=\scriptsize\color{UTPNavy}},
  arr/.style={-{Stealth[length=2mm,width=1.5mm]},
    line width=0.8pt,draw=UTPTeal},
  feedback/.style={arr,draw=UTPMuted,rounded corners=2mm}]
\node[stage] (q) at (0,0) {\textbf{1. Frame}\\[4pt]
  Define the question\\and scope};
\node[stage] (a) at (0.205,0) {\textbf{2. Prepare}\\[4pt]
  Audit quality\\and organise data};
\node[stage] (e) at (0.410,0) {\textbf{3. Explore}\\[4pt]
  Visualise patterns\\and relationships};
\node[stage] (m) at (0.615,0) {\textbf{4. Model}\\[4pt]
  Fit models and\\assess uncertainty};
\node[stage,draw=UTPTeal!80,fill=UTPPaleTeal] (v) at (0.820,0)
  {\textbf{5. Validate}\\[4pt]
   Check results\\and communicate};
\draw[arr] (q.east) -- (a.west);
\draw[arr] (a.east) -- (e.west);
\draw[arr] (e.east) -- (m.west);
\draw[arr] (m.east) -- (v.west);
\draw[feedback] (v.south) -- (0.820,-1.30) --
  node[midway,below=3pt,inner sep=1pt,font=\scriptsize\color{UTPMuted}]
    {New evidence refines the question}
  (0,-1.30) -- (q.south);
\end{tikzpicture}
\end{utpfig}

\begin{utptable}
\begin{tabularx}{\linewidth}{@{}>{\bfseries\leavevmode\color{UTPNavy}\centering\arraybackslash}m{0.045\linewidth}
  D{0.23\linewidth}
  D{0.235\linewidth}
  Y@{}}
\toprule
\rowcolor{UTPPaleBlue!92}
\utpth{Lab} & \utpth{Analytical focus} & \utpth{Dataset(s)} & \utpth{Key outputs} \\
\midrule
1 & Workflow and reproducibility & Bike Sharing (daily) & Data profile, analytical question, and time-aware summaries \\
2 & Data quality & Iris and Titanic & Missingness audit, category checks, and cleaning decisions \\
3 & Wrangling and features & Bike Sharing (hourly) & Parsed dates, derived features, and temporal summaries \\
4 & Exploratory analysis & Medical Costs & Distributions, outliers, subgroup comparisons, and relationships \\
5 & Visual communication & Red Wine Quality & Purposeful graphics with evidence-based interpretations \\
6 & Sampling and inference & Student Performance & Bootstrap interval, group comparison, effect size, and limitations \\
7 & Regression & Ames Housing & Held-out MAE and RMSE, with residual diagnostics \\
8 & Classification & Heart Failure & Predicted probabilities, confusion matrix, sensitivity, specificity, and precision \\
9 & Validation and preprocessing & Credit-card Fraud & Stratified cross-validation and threshold trade-offs \\
10 & Unsupervised learning & Mall Customers and Wine Recognition & Standardised k-means segments and PCA representation \\
\bottomrule
\end{tabularx}
\end{utptable}

\begin{keyidea}
Every practical submission should answer four questions: \emph{What was the question? What data were used? What evidence was produced? What are the limitations?}
\end{keyidea}
